
mais tels qu'on ait
@@M0@@

\textbf{5.3.} Lorsque @@M1@@ est lisse et quasi-projectif sur un corps @@M2@@, on peut de même se poser la question de l'existence d'un homomorphisme d'anneaux
@@M3@@
à valeurs dans l'anneau de Chow @@M4@@, transformant classes de Chern en classes de Chern, et on peut répéter à ce propos les réflexions qui précèdent, qui montrent en particulier comment construire directement @@M5@@ à défaut de @@M6@@ lui-même.

\textbf{5.4.} Signalons, pour terminer ce numéro, qu'essentiellement la même question se pose si, au lieu d'un schéma @@M7@@, on considère un espace topologique compact (disons) @@M8@@, ce qui permet encore de considérer le gradué @@M9@@ de @@M10@@ muni de sa @@M11@@-filtration, et de se proposer de trouver un homomorphisme d'anneaux
@@M12@@
transformant classes de Chern en classes de Chern.

\section*{6. Théorème de Riemann--Roch cohomologique, et homomorphisme de Gysin}

\textbf{6.1.} Pour simplifier, nous nous bornerons ici aux schémas @@M13@@ admettant un Module inversible ample, ce qui implique que l'on n'a pas à distinguer entre le @@M14@@ naif et l'autre. Commençons par la formule de Riemann--Roch ``avec dénominateurs'', pour un morphisme propre @@M15@@ d'intersection complète. Comme (pour un nombre premier @@M16@@ premier aux caractéristiques résiduelles de @@M17@@) on peut définir les classes de Chern @@M18@@-adiques d'un Module localement libre @@M19@@,
@@M20@@
(SGA 5 VII), correspondant à un homomorphisme
@@M21@@
la question se pose de donner une variante cohomologique de la formule de Riemann--Roch. Le seul ingrédient qui manque pour donner un sens à la formule est un \uline{homomorphisme de Gysin}
